\begin{table}[htbp]
\centering
\small
\caption{Bloco 2 (google/gemini-2.5-flash-lite): experimento de diluição.}
\label{tab:bloco2_E1_diluicao_flashlite}
\begin{tabular}{rrrrrr}
\toprule
n & difíceis & fáceis & Acur.\ vs espec. & $\kappa$ & F1-Macro \\
\midrule
4 & 4 & 0 & 58,4 (38,4–81,5) & 0,167 & 58,4 \\
6 & 4 & 2 & 89,8 (85,4–93,8) & 0,796 & 89,8 \\
8 & 4 & 4 & 90,1 (86,6–92,3) & 0,803 & 90,1 \\
14 & 4 & 10 & 89,2 (86,0–91,9) & 0,785 & 89,2 \\
20 & 4 & 16 & 89,7 (86,2–93,8) & 0,795 & 89,7 \\
24 & 4 & 20 & 89,4 (86,5–92,9) & 0,789 & 89,4 \\
\bottomrule
\multicolumn{6}{p{0.95\linewidth}}{\footnotesize Especialista E1. Exemplos difíceis fixos (menor margem) e fáceis acrescentados progressivamente. Acur.\ vs espec.\ = acurácia (LLM vs especialista): fração dos pontos de teste em que o LLM reproduz o rótulo do especialista. $\kappa$ = Cohen's $\kappa$ (faixas de Landis e Koch, 1977). F1-Macro em \%. Médias das sementes (mín–máx).} \\
\end{tabular}
\end{table}
